\section*{Capítulo \ref{ch:b1:units-dimensions} --- Unidades, dimensões e medida}

\begin{solution}{exo:b1:units-dimensions:1}
$\unit{J} = \unit{kg.m^2.s^{-2}}$; $\unit{W} = \unit{kg.m^2.s^{-3}}$;
$\unit{Pa} = \unit{kg.m^{-1}.s^{-2}}$; $\unit{C} = \unit{A.s}$.
$\unit{Pa.m^3} = \unit{kg.m^{-1}.s^{-2}.m^3} = \unit{kg.m^2.s^{-2}} = \unit{J}$
(o trabalho das forças de pressão, $P\,\dd V$).
\end{solution}

\begin{solution}{exo:b1:units-dimensions:2}
Pressão $M L^{-1} T^{-2}$; potência $M L^2 T^{-3}$; carga $I\,T$;
tensão $= $ potência/corrente $= M L^2 T^{-3} I^{-1}$.
$[G] = [F r^2/m^2] = M L T^{-2} L^2 M^{-2} = L^3 M^{-1} T^{-2}$;
$[h] = [E/\nu] = M L^2 T^{-2} \cdot T = M L^2 T^{-1}$;
$[k_B] = [E/T] = M L^2 T^{-2} \Theta^{-1}$.
\end{solution}

\begin{solution}{exo:b1:units-dimensions:3}
$\sqrt{2gh}$: $\sqrt{L T^{-2} \cdot L} = L T^{-1}$, homogênea.
$\tfrac12 mv$: $M L T^{-1} \neq M L^2 T^{-2}$, não é uma energia.
$\rho g h$: $M L^{-3} \cdot L T^{-2} \cdot L = M L^{-1} T^{-2}$, uma pressão.
$v_0 t + \tfrac12 g t^2$: as duas parcelas em $L$, homogênea.
$2\pi\sqrt{g/\ell}$: $\sqrt{T^{-2}} = T^{-1}$ --- uma frequência, não um período.
$I_0\eu^{-t/RC}$: $RC$ em $\unit{\ohm.F} = \unit{s}$, argumento \omterm{def:b1:units-dimensions:dimension}{adimensional}, homogênea.
\end{solution}

\begin{solution}{exo:b1:units-dimensions:4}
Régua: $\delta = \qty{1}{mm}$, $u = 1/\sqrt{12} = \qty{0.29}{mm}$.
Balança: $u = 0.01/\sqrt{12} = \qty{0.003}{g}$.
Voltímetro: $a = 0.005 \times 4.80 = \qty{0.024}{V}$,
$u = 0.024/\sqrt3 = \qty{0.014}{V}$.
\end{solution}

\begin{solution}{exo:b1:units-dimensions:5}
$v = k F^\alpha \mu^\beta$: $L T^{-1} = (M L T^{-2})^\alpha (M L^{-1})^\beta$,
logo $\alpha + \beta = 0$, $\alpha - \beta = 1$, $-2\alpha = -1$:
$\alpha = \tfrac12$, $\beta = -\tfrac12$, $v = k\sqrt{F/\mu}$ (de fato $k = 1$).
Som: $P/\rho$ tem \omterm{def:b1:units-dimensions:dimension}{dimensão} $M L^{-1} T^{-2} / (M L^{-3}) = L^2 T^{-2}$,
logo $v = k\sqrt{P/\rho}$ ($k = \sqrt\gamma \approx 1.2$ para o ar).
\end{solution}

\begin{solution}{exo:b1:units-dimensions:6}
Soma $= \qty{3620}{ms}$, $\bar t = \qty{452.5}{ms}$. Desvios
$-0.5, -5.5, 8.5, 2.5, -3.5, 5.5, -8.5, 1.5$; os quadrados somam $226$;
$s = \sqrt{226/7} = \qty{5.7}{ms}$; $u(\bar t) = 5.7/\sqrt8 = \qty{2.0}{ms}$.
Resultado: $t = \qty{452.5\pm2.0}{ms}$.
\end{solution}

\begin{solution}{exo:b1:units-dimensions:7}
$R = 4.80/0.0124 = \qty{387}{\ohm}$. Incertezas relativas $0.03/4.80 =
0.63\%$ e $0.2/12.4 = 1.6\%$; combinadas $\sqrt{0.63^2 + 1.6^2} = 1.7\%$,
$u(R) = \qty{7}{\ohm}$: $R = \qty{387\pm7}{\ohm}$. A corrente
domina --- melhorar primeiro o amperímetro (ou a sua escala).
\end{solution}

\begin{solution}{exo:b1:units-dimensions:8}
$z = 7/\sqrt{16 + 9} = 1.4 \leq 2$: compatíveis.
Terceiro contra o primeiro: $z = 9/\sqrt{4 + 16} = 2.0$: no limite, ainda
compatíveis. Terceiro contra o segundo: $z = 16/\sqrt{4 + 9} = 4.4$: incompatíveis
--- pelo menos um dos dois tem um erro sistemático ou uma
incerteza subestimada.
\end{solution}

\begin{solution}{exo:b1:units-dimensions:9}
$V = \pi d^3/6 = \pi \times 25.40^3/6 = \qty{8580}{mm^3}$.
$u(V)/V = 3\,u(d)/d = 3 \times 0.02/25.40 = 0.24\%$,
$u(V) = \qty{20}{mm^3}$: $V = \qty{8.58\pm0.02}{cm^3}$.
\end{solution}

\begin{solution}{exo:b1:units-dimensions:10}
$F = k\eta^a R^b v^c$: em $M$: $a = 1$; em $T$: $-a - c = -2$, $c = 1$;
em $L$: $-a + b + c = 1$, $b = 1$: $F = k\eta R v$ (Stokes, $k = 6\pi$).
Com $\rho$ no lugar: em $M$: $a = 1$; em $T$: $-c = -2$, $c = 2$; em $L$:
$-3 + b + 2 = 1$, $b = 2$: $F = k\rho R^2 v^2$.
Gota de chuva: $\eta R v = 1.8\times10^{-5} \times 10^{-3} \times 5 \approx
\qty{1e-7}{N}$; $\rho R^2 v^2 = 1.2 \times 10^{-6} \times 25 =
\qty{3e-5}{N}$, trezentas vezes maior: aplica-se o regime inercial (de
alta velocidade) --- a razão $\rho v R/\eta \approx 330$ é o número de
Reynolds.
\end{solution}

\begin{solution}{exo:b1:units-dimensions:11}
$\bar I = \qty{5.0}{mA}$, $\bar U = \qty{1.000}{V}$; desvios em $I$:
$-3, -1, 1, 3$; em $U$: $-0.59, -0.22, 0.22, 0.59$.
$\sum (I_i - \bar I)(U_i - \bar U) = 3.98$, $\sum (I_i - \bar I)^2 = 20$:
$a = \qty{0.199}{V/mA} = \qty{199}{\ohm}$, $b = 1.000 - 0.199 \times 5 =
\qty{0.005}{V}$. Com $u(U) = \qty{0.02}{V}$ por ponto, um coeficiente linear de
\qty{0.005}{V} está bem dentro da incerteza: compatível com zero, o
resistor é ôhmico, $R \approx \qty{199}{\ohm}$.
\end{solution}

\begin{solution}{exo:b1:units-dimensions:12}
$\ell = G^a h^b c^c$: em $M$: $-a + b = 0$; em $T$: $-2a - b - c = 0$; em $L$:
$3a + 2b + c = 1$. Logo $b = a$, $c = -3a$, $2a = 1$:
$\ell_P = \sqrt{Gh/c^3} = \sqrt{4.42\times10^{-44}/2.7\times10^{25}}
= \qty{4.1e-35}{m}$; $t_P = \ell_P/c = \qty{1.4e-43}{s}$;
$m_P = \sqrt{hc/G} = \sqrt{1.99\times10^{-25}/6.67\times10^{-11}}
= \qty{5.5e-8}{kg}$. O comprimento de Planck fica vinte \omterm{def:b1:units-dimensions:oom}{ordens de grandeza}
abaixo de um próton; a massa de Planck, \qty{55}{\micro g}, é um grão de poeira
--- enorme para uma partícula, ínfima para uma pessoa ($10^9$ vezes mais leve).
\end{solution}

\begin{solution}{pb:b1:units-dimensions:1}
\textbf{1.} $\sqrt{L/(L T^{-2})} = \sqrt{T^2} = T$.

\textbf{2.} $\tau = k\ell^\alpha m^\beta g^\gamma$ dá
$T = L^{\alpha+\gamma} M^\beta T^{-2\gamma}$, logo $\beta = 0$: nenhuma
combinação de $\ell$ e $g$ pode cancelar uma massa.

\textbf{3.} $\theta_0$ é \omterm{def:b1:units-dimensions:dimension}{adimensional}, de modo que qualquer fator
$f(\theta_0)$ é invisível às dimensões. Mantenha a amplitude pequena
(abaixo de cerca de $10^\circ$), onde $f \approx 1$ com erro menor que $0.2\%$.

\textbf{4.} $g = 4\pi^2\ell/\tau^2$.

\textbf{5.} $\ell = 99.0 + 1.00 = \qty{100.0}{cm} = \qty{1.000}{m}$.

\textbf{6.} $u = \qty{1}{mm}/\sqrt{12} = \qty{0.29}{mm}$.

\textbf{7.} Alinhamento: $5/\sqrt3 = \qty{2.9}{mm}$; combinadas
$\sqrt{0.29^2 + 2.9^2} = \qty{2.9}{mm}$ --- a leitura da trena é
desprezível. $u(\ell)/\ell = 0.29\%$.

\textbf{8.} O erro de tempo de reação (cerca de \qty{0.1}{s}) é o mesmo
quer se cronometre um ou dez períodos, de modo que o seu efeito sobre um
período fica dez vezes menor. Repetir revela a dispersão (daí $s$) e
divide a incerteza da média por $\sqrt8$.

\textbf{9.} Soma $= \qty{160.80}{s}$, $\overline{t_{10}} = \qty{20.100}{s}$.
Desvios $0.02, -0.05, 0.08, -0.01, 0.05, -0.08, 0.01, -0.02$; os quadrados
somam $\num{188e-4}$; $s = \sqrt{188\times10^{-4}/7} = \qty{0.052}{s}$.

\textbf{10.} $u(\overline{t_{10}}) = 0.052/\sqrt8 = \qty{0.018}{s}$;
$\tau = \qty{2.0100}{s}$, $u(\tau) = \qty{0.0018}{s}$.

\textbf{11.} $u(\tau)/\tau = 0.09\%$, três vezes menor que
$u(\ell)/\ell = 0.29\%$.

\textbf{12.} $g = 4\pi^2 \times 1.000/2.0100^2 = 39.48/4.040 =
\qty{9.772}{m/s^2}$.

\textbf{13.} Como $g \propto \ell\,\tau^{-2}$:
$\dfrac{u(g)}{g} = \sqrt{\left(\dfrac{u(\ell)}{\ell}\right)^2
+ \left(2\dfrac{u(\tau)}{\tau}\right)^2} = \sqrt{0.29^2 + 0.18^2}\,\%
= 0.34\%$.

\textbf{14.} $u(g) = 0.0034 \times 9.772 = \qty{0.033}{m/s^2}$:
$g = \qty{9.77\pm0.03}{m/s^2}$.

\textbf{15.} $z = \abs{9.772 - 9.81}/0.033 = 1.2 \leq 2$: compatíveis.

\textbf{16.} O comprimento ($0.29\%$ contra $0.18\%$ para o termo do período).
Localizar as duas extremidades a $\pm\qty{1}{mm}$ leva $u(\ell)/\ell$ a
$0.06\%$ e $u(g)/g$ a $\sqrt{0.06^2 + 0.18^2} = 0.19\%$; dobrar
o comprimento reduz $u(\ell)/\ell$ à metade e alonga $\tau$, ajudando os
dois termos. Cronometrar vinte períodos em vez de dez reduz à metade só o
termo do período: $u(g)/g = \sqrt{0.29^2 + 0.09^2} = 0.30\%$ --- ganho
pequeno enquanto o comprimento é o gargalo.

\textbf{17.} Não: um deslocamento comum a todas as leituras é um
\emph{erro sistemático} (um viés); o orçamento cobre a dispersão aleatória
e as tolerâncias declaradas, não um desvio constante. Ele é detectado
mudando de método ou de instrumento, ou --- Parte V --- por um
coeficiente linear não nulo da reta $\tau^2$--$\ell$.

\textbf{18.} $\tau^2 = (4\pi^2/g)\,\ell$ é uma reta que passa pela
origem: a linearidade e o coeficiente linear nulo são ambos testáveis a
olho e por \omterm{prop:b1:units-dimensions:regression}{mínimos quadrados}; $\tau \propto \sqrt\ell$ não é.

\textbf{19.} $\bar\ell = \qty{0.800}{m}$; $\overline{\tau^2} = 16.16/5 =
\qty{3.232}{s^2}$.

\textbf{20.} Desvios em $\ell$: $-0.4, -0.2, 0, 0.2, 0.4$; em
$\tau^2$: $-1.612, -0.822, 0.018, 0.808, 1.608$. Soma cruzada $= 1.614$,
$\sum(\ell_i - \bar\ell)^2 = 0.400$: $a = \qty{4.035}{s^2/m}$,
$b = 3.232 - 4.035 \times 0.800 = \qty{0.004}{s^2}$.

\textbf{21.} $g = 4\pi^2/a = 39.48/4.035 = \qty{9.78}{m/s^2}$.

\textbf{22.} $b \approx 0$, como o modelo prevê. Um $b$ claramente não
nulo denunciaria um desvio sistemático $\delta$ em todos os comprimentos
($\tau^2 = a(\ell_{\text{med}} - \delta)$ dá $b = -a\delta$), por exemplo\ um
centro de esfera mal localizado --- ou um efeito de amplitude.

\textbf{23.} $u(g)/g = u(a)/a = 0.05/4.035 = 1.2\%$, $u(g) =
\qty{0.12}{m/s^2}$: $g = \qty{9.78\pm0.12}{m/s^2}$.

\textbf{24.} $z = \abs{9.772 - 9.784}/\sqrt{0.033^2 + 0.12^2} = 0.1$:
plenamente compatíveis. A medida única e cuidadosa é mais precisa; a
reta valida o modelo e protege contra um desvio de comprimento.

\textbf{25.} $\tau = 2\pi\sqrt{10/9.8} = \qty{6.3}{s}$; uma única
oscilação cronometrada com tempo de reação de cerca de \qty{0.2}{s} dá
$u(\tau)/\tau \approx
3\%$, logo $u(g)/g \approx 6\%$, enquanto
$u(\ell)/\ell$ cai para $0.03\%$: vinte vezes pior que a Parte IV. O ganho
de comprimento é desperdiçado ao cronometrar um único período. Melhor
estratégia: pêndulo longo \emph{e} muitos períodos, repetidos --- dez
oscilações do pêndulo de \qty{10}{m}, oito vezes, dão
$u(\tau)/\tau \approx 0.03\%$ e $u(g)/g \approx 0.07\%$.
\end{solution}
